\documentclass{article}
\usepackage[utf8]{inputenc}
\usepackage{geometry}
\geometry{
  a4paper,
  total={170mm,257mm},
  left=20mm,
  top=20mm,
}
\usepackage{graphicx}
\usepackage{titling}
\usepackage{booktabs}
\usepackage{listings}
\usepackage{xcolor}
\usepackage{hyperref}
\usepackage{fancyvrb}
\usepackage{float}
\usepackage{amsmath}

\definecolor{darkblue}{rgb}{0.0, 0.0, 0.5}
\RecustomVerbatimEnvironment{verbatim}{Verbatim}{formatcom=\color{darkblue}}
\setlength{\parindent}{0pt}

\title{\textbf{DEEP LEARNING} \\
{\Large \textbf{Report 3: Logistic Regression}}}
\date{May 2026}

\usepackage{fancyhdr}
\fancypagestyle{plain}{
  \fancyhf{}
  \fancyfoot[L]{\thedate}
  \fancyfoot[R]{\thepage}
  \fancyhead[L]{University of Science and Technology of Hanoi}
  \fancyhead[R]{Deep Learning}
}
\makeatletter
\def\@maketitle{%
  \newpage
  \vskip 1em
  \begin{center}
  \let \footnote \thanks
    {\LARGE \@title \par}%
    \vskip 1em
  \end{center}
  \par
  \vskip 1em}
\makeatother

\begin{document}
\pagestyle{plain}
\maketitle

\noindent\begin{tabular}{@{}ll}
  Student: & Vu Ha Vy -- 2540112 \\
\end{tabular}

\section*{Implementation}

\textbf{Loss function}

Unlike linear regression, logistic regression predicts a probability using
the sigmoid function $\sigma(z) = \frac{1}{1 + e^{-z}}$, where
$z = w_1 x_1 + w_2 x_2 + w_0$.
The model is trained by maximizing the log-likelihood, which is equivalent
to minimizing the binary cross-entropy loss:
\[
  J(w_0, w_1, w_2) = -\frac{1}{N} \sum_{i=1}^{N}
  \Bigl[ y_i \cdot z_i - \log(1 + e^{z_i}) \Bigr]
\]

\begin{verbatim}
def loss(data, w0, w1, w2):
    N = len(data)
    total = 0.0

    for x1, x2, yi in data:
        y_hat = w1 * x1 + w2 * x2 + w0
        z = yi * y_hat - math.log(1 + math.exp(y_hat))
        total += z

    return -(1 / N) * total
\end{verbatim}

\textbf{Gradient descent}

The weights are updated iteratively using the partial derivatives of $J$.
The error at each point is $\hat{y}_i - y_i$, where
$\hat{y}_i = \sigma(z_i)$ is the predicted probability:
\[
  \frac{\partial J}{\partial w_0} = \frac{1}{N}\sum_{i=1}^{N}(\hat{y}_i - y_i)
  \qquad
  \frac{\partial J}{\partial w_1} = \frac{1}{N}\sum_{i=1}^{N} x_{1i}(\hat{y}_i - y_i)
  \qquad
  \frac{\partial J}{\partial w_2} = \frac{1}{N}\sum_{i=1}^{N} x_{2i}(\hat{y}_i - y_i)
\]

\begin{verbatim}
def gradient_descent(data, lr, thres):
    w0 = 0.0
    w1 = 0.0
    w2 = 0.0

    N = len(data)
    i = 0
    J_prev = float('inf')

    while True:
        i += 1

        J = loss(data, w0, w1, w2)

        if abs(J - J_prev) < thres:
            print(f"Converged at iteration {i}, J = {J}")
            break

        J_prev = J

        print(f"{i} | {w0:.4f} | {w1:.4f} | {w2:.4f} | {J:.4f}")

        dw0 = 0.0
        dw1 = 0.0
        dw2 = 0.0

        for x1, x2, yi in data:
            z = w1 * x1 + w2 * x2 + w0
            prediction = sigmoid(z)
            error = prediction - yi

            dw0 += error
            dw1 += x1 * error
            dw2 += x2 * error

        dw0 /= N
        dw1 /= N
        dw2 /= N

        w0 = w0 - lr * dw0
        w1 = w1 - lr * dw1
        w2 = w2 - lr * dw2

    return w0, w1, w2
\end{verbatim}

\section*{Dataset}

The dataset \texttt{loan2.csv} contains 19 data points with two features,
Experience (years) and Salary (million), and a binary label indicating
whether the loan was approved (1) or refused (0).

\begin{table}[H]
\centering
\begin{tabular}{ccc}
\toprule
Experience (year) & Salary (million) & Loan \\
\midrule
3.00  & 4  & 1 \\
2.50  & 4  & 1 \\
1.00  & 4  & 0 \\
2.50  & 5  & 1 \\
2.00  & 5  & 1 \\
1.50  & 5  & 0 \\
0.50  & 5  & 0 \\
1.75  & 6  & 1 \\
0.25  & 6  & 0 \\
1.00  & 7  & 1 \\
0.25  & 7  & 0 \\
0.20  & 7  & 0 \\
0.15  & 7  & 0 \\
2.00  & 8  & 1 \\
1.00  & 8  & 0 \\
0.15  & 8  & 0 \\
0.10  & 8  & 0 \\
0.50  & 9  & 1 \\
1.00  & 10 & 1 \\
\bottomrule
\end{tabular}
\caption{Contents of \texttt{loan2.csv}}
\end{table}

\section*{Results}

With learning rate $r = 0.1$ and threshold $t = 10^{-6}$, the algorithm
converged after 13277 iterations with final loss $J = 0.2058$.
The learned model is:
\[
  \hat{y} = \sigma(-13.5672 + 4.5400 \cdot x_1 + 0.2058 \cdot x_2)
\]
where $x_1$ is Experience, $x_2$ is Salary, and
$\sigma(z) = \frac{1}{1+e^{-z}}$.

\begin{figure}[H]
    \centering
    \includegraphics[width=0.6\textwidth]{plot.png}
    \caption{Decision boundary (green) separating approved loans (red) from
    refused loans (blue), plotted over Salary vs.\ Experience.}
\end{figure}

\textbf{Effects of different learning rates}

Table below compares three learning rates with the same threshold
$t = 10^{-6}$. A larger $r$ converges faster and also reaches a lower loss,
because bigger steps escape flat regions of the loss surface more quickly.
However, setting $r$ too large risks overshooting and diverging.
Here $r = 0.1$ gives the best result, converging in 13277 iterations with
the lowest loss of 0.2058.

\begin{table}[H]
    \centering
    \begin{tabular}{cccccc}
    \toprule
    Learning rate $r$ & Iterations & Final loss $J$ & $w_0$ & $w_1$ & $w_2$ \\
    \midrule
    $0.001$ & 76,147 & 0.3460 & $-3.1537$ & $2.1585$ & $0.1413$ \\
    $0.01$  & 41,356 & 0.2359 & $-8.7517$ & $3.3970$ & $0.7597$ \\
    $0.1$   & 13,277 & 0.2058 & $-13.5672$ & $4.5400$ & $0.2058$ \\
    \bottomrule
    \end{tabular}
    \caption{Comparison of learning rates (threshold $t = 10^{-6}$)}
    \label{tab:lr}
\end{table}

\end{document}